\documentclass[a4paper,12pt]{article}

% ==============================================
% CẤU HÌNH CHO PDFLATEX (NHANH & KHÔNG TIMEOUT)
% ==============================================
\usepackage[utf8]{inputenc}
\usepackage{graphicx}
\usepackage[T5]{fontenc}
\usepackage[vietnamese]{babel}
\usepackage{mathptmx} % Font Times New Roman giả lập
\usepackage{geometry}
\usepackage{graphicx}
\usepackage{longtable}
\usepackage{listings}
\usepackage{xcolor} % nếu muốn highlight
\usepackage{graphicx} % Để chèn ảnh
\usepackage{float}    % Để ép ảnh nằm đúng chỗ (quan trọng)
\usepackage{float}
\usepackage{hyperref}
\usepackage{listings}
\usepackage{xcolor}

\geometry{left=2.5cm, right=2cm, top=2cm, bottom=2cm}

% --- CẤU HÌNH HIỂN THỊ CODE TIẾNG VIỆT ---
\definecolor{codegreen}{rgb}{0,0.6,0}
\definecolor{codegray}{rgb}{0.5,0.5,0.5}
\definecolor{codepurple}{rgb}{0.58,0,0.82}
\definecolor{backcolour}{rgb}{0.95,0.95,0.92}

\lstdefinestyle{mystyle}{
    backgroundcolor=\color{backcolour},   
    commentstyle=\color{codegreen},
    keywordstyle=\color{blue},
    numberstyle=\tiny\color{codegray},
    stringstyle=\color{codepurple},
    basicstyle=\ttfamily\footnotesize,
    breakatwhitespace=false,         
    breaklines=true,                 
    captionpos=b,                    
    keepspaces=true,                 
    numbers=left,                    
    numbersep=5pt,                  
    showspaces=false,                
    showstringspaces=false,
    showtabs=false,                  
    tabsize=2,
    frame=single,
    % PHẦN QUAN TRỌNG NHẤT: Bảng thay thế ký tự tiếng Việt
    literate={đ}{{\dj}}1 {Đ}{{\DJ}}1
    {á}{{\'a}}1 {à}{{\`a}}1 {ả}{{\h a}}1 {ã}{{\~a}}1 {ạ}{{\d a}}1
    {ă}{{\u a}}1 {ắ}{{\'a}}1 {ằ}{{\`a}}1 {ẳ}{{\h a}}1 {ẵ}{{\~a}}1 {ặ}{{\d a}}1
    {â}{{\^a}}1 {ấ}{{\'a}}1 {ầ}{{\`a}}1 {ẩ}{{\h a}}1 {ẫ}{{\~a}}1 {ậ}{{\d a}}1
    {é}{{\'e}}1 {è}{{\`e}}1 {ẻ}{{\h e}}1 {ẽ}{{\~e}}1 {ẹ}{{\d e}}1
    {ê}{{\^e}}1 {ế}{{\'e}}1 {ề}{{\`e}}1 {ể}{{\h e}}1 {ễ}{{\~e}}1 {ệ}{{\d e}}1
    {í}{{\'i}}1 {ì}{{\`i}}1 {ỉ}{{\h i}}1 {ĩ}{{\~i}}1 {ị}{{\d i}}1
    {ó}{{\'o}}1 {ò}{{\`o}}1 {ỏ}{{\h o}}1 {õ}{{\~o}}1 {ọ}{{\d o}}1
    {ô}{{\^o}}1 {ố}{{\'o}}1 {ồ}{{\`o}}1 {ổ}{{\h o}}1 {ỗ}{{\~o}}1 {ộ}{{\d o}}1
    {ơ}{{\ur o}}1 {ớ}{{\'o}}1 {ờ}{{\`o}}1 {ở}{{\h o}}1 {ỡ}{{\~o}}1 {ợ}{{\d o}}1
    {ú}{{\'u}}1 {ù}{{\`u}}1 {ủ}{{\h u}}1 {ũ}{{\~u}}1 {ụ}{{\d u}}1
    {ư}{{\ur u}}1 {ứ}{{\'u}}1 {ừ}{{\`u}}1 {ử}{{\h u}}1 {ữ}{{\~u}}1 {ự}{{\d u}}1
    {ý}{{\'y}}1 {ỳ}{{\`y}}1 {ỷ}{{\h y}}1 {ỹ}{{\~y}}1 {ỵ}{{\d y}}1
}

\lstdefinelanguage{JavaScript}{
  keywords={typeof, new, true, false, catch, function, return, null, switch, var, if, while, do, else, case, break},
  ndkeywords={class, export, boolean, throw, implements, import, this},
  sensitive=true,
  comment=[l]{//},
  morecomment=[s]{/*}{*/},
  morestring=[b]',
  morestring=[b]"
}


\lstset{style=mystyle}

\begin{document}

% ==============================================
% 2. TRANG BÌA
% ==============================================
\begin{titlepage}
    \centering
    \textbf{TRƯỜNG ĐẠI HỌC SÀI GÒN}\\
    \textbf{KHOA CÔNG NGHỆ THÔNG TIN}
    
    \vspace{2cm}
    % \includegraphics[width=1.35in]{media/image4.jpg} % Placeholder
    
    {\Huge \textbf{KIỂM THỬ PHẦN MỀM}}
    
    \vspace{1cm}
    % \includegraphics[width=6in,keepaspectratio]{media/image82.png} % Placeholder
    
    \vspace{1cm}
    \begin{quote}
    \textbf{ĐỀ TÀI :}
    
    \textbf{Ứng dụng Đăng nhập \& Quản lý Sản phẩm}
    \end{quote}
    
    \subsection*{(Login \& Product Management - Version 1.0)}
    
    \vspace{2cm}
    
    \begin{flushleft}
    \textbf{GIẢNG VIÊN:} Từ Lãng Phiêu \\
    \textbf{NHÓM:} 04 \\
    \textbf{SINH VIÊN:} \\
    Kiều Hoài Nam - 3123410226 \\
    Lê Quốc Cường - 3123410040 \\
    Lê Công Vinh - 3123410431 \\
    Phan Vĩnh Thái - 3123560079 \\
    Nguyễn Dương Khang - 3123410151
    \end{flushleft}
    
    \vfill
    \centering
    \textbf{TP. HỒ CHÍ MINH, THÁNG 11/2025}
\end{titlepage}

\newpage
\tableofcontents
\newpage

% ==============================================
% 3. NỘI DUNG CHI TIẾT
% ==============================================

\section*{Đóng Góp Của Nhóm}
\addcontentsline{toc}{section}{Đóng Góp Của Nhóm}

\begin{longtable}{|p{0.25\linewidth}|p{0.2\linewidth}|p{0.15\linewidth}|p{0.3\linewidth}|}
\hline
\textbf{Thành viên} & \textbf{Chức vụ} & \textbf{Phần trăm} & \textbf{Đóng góp} \\
\hline
\endhead
Lê Quốc Cường & Trưởng nhóm & & \\ \hline
Kiều Hoài Nam & Thành viên & & \\ \hline
Lê Công Vinh & Thành viên & & \\ \hline
Phan Vĩnh Thái & Thành viên & & \\ \hline
Nguyễn Dương Khang & Thành viên & & \\ \hline
\end{longtable}

\section{GIỚI THIỆU DỰ ÁN}

\subsection{Tổng quan}
Dự án \textbf{FloginFE\_BE} là một ứng dụng web full-stack được phát triển để thực hành các kỹ thuật kiểm thử phần mềm. Hệ thống bao gồm các chức năng chính:
\begin{itemize}
    \item \textbf{Login:} Hệ thống đăng nhập bảo mật với validation chặt chẽ.
    \item \textbf{Product Management:} Quản lý sản phẩm với đầy đủ các thao tác CRUD (Create, Read, Update, Delete).
\end{itemize}

\subsection{Công nghệ sử dụng}
\begin{itemize}
    \item \textbf{Frontend:} React 18+, Jest, React Testing Library.
    \item \textbf{Backend:} Spring Boot 3.2+, JUnit 5, Mockito.
    \item \textbf{Database:} MySQL.
    \item \textbf{CI/CD:} GitHub Actions.
\end{itemize}

\section{CÂU 1: PHÂN TÍCH VÀ THIẾT KẾ TEST CASES}

\subsection{Câu 1.1: Login - Phân tích và Test Scenarios (10 điểm)}

\subsubsection{Yêu cầu (5 điểm)}
\textbf{a) Phân tích yêu cầu chức năng Login (2 điểm)}

\textbf{1. Validation rules cho Username}
\begin{itemize}
    \item Bắt buộc nhập (required, không được để trống).
    \item Độ dài hợp lệ: 3--50 ký tự.
    \item Chỉ cho phép ký tự: a--z, A--Z, 0--9, dấu gạch ngang (-), dấu chấm (.), dấu gạch dưới (\_).
    \item Không được chứa khoảng trắng đầu/cuối.
    \item Không được chứa ký tự đặc biệt ngoài danh sách cho phép.
\end{itemize}

\textbf{2. Validation rules cho Password}
\begin{itemize}
    \item Bắt buộc nhập.
    \item Độ dài hợp lệ: 6--100 ký tự.
    \item Phải chứa ít nhất 1 chữ cái (a--z hoặc A--Z).
    \item Phải chứa ít nhất 1 chữ số (0--9).
    \item Không giới hạn ký tự đặc biệt (nếu hệ thống không cấm).
    \item Không có khoảng trắng đầu/cuối (trim).
\end{itemize}

\textbf{3. Authentication Flow (Luồng xác thực)}
\begin{enumerate}
    \item User nhập username + password.
    \item Hệ thống validate input: Nếu invalid $\rightarrow$ trả về lỗi Validation và không gọi API backend.
    \item Nếu input hợp lệ $\rightarrow$ gửi request đến backend: POST /login \{ username, password \}.
    \item Backend kiểm tra: Có tồn tại username không? Password có khớp không?
    \item Nếu hợp lệ: Trả về token (JWT), Frontend lưu token, Chuyển hướng người dùng sang products.
    \item Nếu không hợp lệ: Trả về lỗi: Tên đăng nhập hoặc mật khẩu không đúng.
\end{enumerate}

\textbf{4. Error Handling}
\begin{itemize}
    \item Username rỗng $\rightarrow$ "Username không được để trống".
    \item Password rỗng $\rightarrow$ "Password không được để trống".
    \item Username độ dài nhỏ hơn 3 hoặc lớn hơn 50 $\rightarrow$ "Username phải từ 3-50 ký tự".
    \item Password độ dài nhỏ hơn 6 hoặc lớn hơn 100 $\rightarrow$ "Password phải từ 6-100 ký tự".
    \item Username sai format $\rightarrow$ "Username chỉ được chứa chữ cái, số...".
    \item Password sai format $\rightarrow$ "Password phải chứa ít nhất 1 chữ cái và 1 chữ số".
    \item Username không tồn tại $\rightarrow$ "User không tồn tại!".
    \item Sai mật khẩu hoặc username $\rightarrow$ "Tên đăng nhập hoặc mật khẩu không đúng!".
\end{itemize}

\subsubsection{Test Scenarios cho Login (2 điểm)}
\begin{longtable}{|p{0.05\linewidth}|p{0.6\linewidth}|p{0.25\linewidth}|}
\hline
\textbf{STT} & \textbf{Test Scenario} & \textbf{Loại test} \\
\hline
1 & Đăng nhập thành công với username + password hợp lệ & Happy Path \\ \hline
2 & Username rỗng & Negative \\ \hline
3 & Password rỗng & Negative \\ \hline
4 & Username \textless{} 3 ký tự & Boundary -- Invalid \\ \hline
5 & Username \textgreater{} 50 ký tự & Boundary -- Invalid \\ \hline
6 & Password \textless{} 6 ký tự & Boundary -- Invalid \\ \hline
7 & Password \textgreater{} 100 ký tự & Boundary -- Invalid \\ \hline
8 & Username chứa ký tự không hợp lệ & Edge case \\ \hline
9 & Password không chứa số & Negative (Format) \\ \hline
10 & Password không chứa chữ cái & Negative (Format) \\ \hline
11 & Sai username & Negative \\ \hline
12 & Sai password & Negative \\ \hline
13 & Username có khoảng trắng đầu/cuối & Edge case \\ \hline
14 & Password có khoảng trắng đầu/cuối & Edge case \\ \hline
15 & Server trả lỗi 500 & Error handling \\ \hline
\end{longtable}
\textbf{Phân loại ưu tiên test scenarios (1 điểm)}
\begin{itemize}
    \item \textbf{Critical}
    \begin{itemize}
        \item Đăng nhập thành công (TC1) $\rightarrow$ Chức năng chính, ảnh hưởng toàn bộ hệ thống.
        \item Sai username hoặc password (TC11, TC12) $\rightarrow$ Ảnh hưởng trực tiếp đến bảo mật.
        \item Username/password rỗng (TC2, TC3) $\rightarrow$ Validation cơ bản bắt buộc hoạt động.
    \end{itemize}
    \item \textbf{High}
    \begin{itemize}
        \item Password không chứa số/chữ (TC9, TC10) $\rightarrow$ Vi phạm quy định bảo mật.
        \item Username chứa ký tự không hợp lệ (TC8).
    \end{itemize}
    \item \textbf{Medium}
    \begin{itemize}
        \item Boundary tests (TC4, TC5, TC6, TC7) $\rightarrow$ Đảm bảo tính ổn định input.
    \end{itemize}
    \item \textbf{Low}
    \begin{itemize}
        \item Khoảng trắng đầu cuối (TC13, TC14).
        \item Server 500 (TC15) $\rightarrow$ Không xảy ra thường xuyên nhưng cần kiểm tra.
    \end{itemize}
\end{itemize}

\subsubsection{Thiết kế Test Cases chi tiết (5 điểm)}

% Test Case 1
\begin{longtable}{|p{0.3\linewidth}|p{0.65\linewidth}|}
\hline
\textbf{Test Case ID} & TC\_LOGIN\_001 \\ \hline
\textbf{Test Name} & Đăng nhập thành công với credentials hợp lệ \\ \hline
\textbf{Priority} & Critical \\ \hline
\textbf{Preconditions} & - User account đã tồn tại \newline - Application đang chạy \\ \hline
\textbf{Test Steps} & 1. Navigate to login page \newline 2. Enter valid username \newline 3. Enter valid password \newline 4. Click Login button \\ \hline
\textbf{Test Data} & Username: testuser \newline Password: Test123 \\ \hline
\textbf{Expected Result} & - Success message displayed \newline - Token stored \newline - Redirect to dashboard \\ \hline
\textbf{Status} & Not Run \\ \hline
\end{longtable}

% Test Case 2
\begin{longtable}{|p{0.3\linewidth}|p{0.65\linewidth}|}
\hline
\textbf{Test Case ID} & TC\_LOGIN\_002 \\ \hline
\textbf{Test Name} & Đăng nhập thất bại với password sai \\ \hline
\textbf{Priority} & High \\ \hline
\textbf{Preconditions} & - User account đã tồn tại \newline - Application đang chạy \\ \hline
\textbf{Test Steps} & 1. Navigate to login page \newline 2. Enter valid username \newline 3. Enter invalid password \newline 4. Click Login button \\ \hline
\textbf{Test Data} & Username: testuser \newline Password: WrongPass123 \\ \hline
\textbf{Expected Result} & - Error message hiển thị: "Tên đăng nhập hoặc mật khẩu không đúng!" \newline - Không lưu token \newline - Vẫn ở trang login \\ \hline
\textbf{Status} & Not Run \\ \hline
\end{longtable}

% Test Case 3
\begin{longtable}{|p{0.3\linewidth}|p{0.65\linewidth}|}
\hline
\textbf{Test Case ID} & TC\_LOGIN\_003 \\ \hline
\textbf{Test Name} & Đăng nhập thất bại với username không tồn tại \\ \hline
\textbf{Priority} & High \\ \hline
\textbf{Test Steps} & 1. Navigate to login page \newline 2. Enter invalid username \newline 3. Enter any password \newline 4. Click Login button \\ \hline
\textbf{Test Data} & Username: nonexistentuser \newline Password: AnyPassword123 \\ \hline
\textbf{Expected Result} & - Error message hiển thị: "Đăng nhập thất bại!" \newline - Không lưu token \newline - Vẫn ở trang login \\ \hline
\textbf{Status} & Not Run \\ \hline
\end{longtable}

% Test Case 4
\begin{longtable}{|p{0.3\linewidth}|p{0.65\linewidth}|}
\hline
\textbf{Test Case ID} & TC\_LOGIN\_004 \\ \hline
\textbf{Test Name} & Đăng nhập thất bại khi để trống username/password \\ \hline
\textbf{Priority} & High \\ \hline
\textbf{Test Steps} & 1. Navigate to login page \newline 2. Để trống username field \newline 3. Để trống password field \newline 4. Click Login button \\ \hline
\textbf{Expected Result} & - Validation message hiển thị: "Username không được để trống" và "Password không được để trống" \\ \hline
\textbf{Status} & Not Run \\ \hline
\end{longtable}

% Test Case 5
\begin{longtable}{|p{0.3\linewidth}|p{0.65\linewidth}|}
\hline
\textbf{Test Case ID} & TC\_LOGIN\_005 \\ \hline
\textbf{Test Name} & Đăng nhập thất bại với password không đủ độ dài \\ \hline
\textbf{Priority} & Medium \\ \hline
\textbf{Test Steps} & 1. Navigate to login page \newline 2. Enter valid username \newline 3. Enter short password \newline 4. Click Login button \\ \hline
\textbf{Test Data} & Username: testuser \newline Password: pass \\ \hline
\textbf{Expected Result} & - Validation message hiển thị: "Password phải từ 6-100 ký tự" \newline - Không gửi request đến server \\ \hline
\textbf{Status} & Not Run \\ \hline
\end{longtable}

\subsection{Câu 1.2: Product - Phân tích và Test Scenarios (10 điểm)}

\subsubsection{Yêu cầu (5 điểm)}
\textbf{Phân tích yêu cầu chức năng Product CRUD (2 điểm)}

\textbf{1. Create -- Thêm sản phẩm mới}
\begin{itemize}
    \item name: không rỗng, tối đa 200 ký tự.
    \item description: tối đa 1000 ký tự.
    \item price: bắt buộc, \textgreater{} 0.
    \item quantity: \textgreater{}= 0.
    \item category: tối đa 50 ký tự, phải thuộc danh sách hợp lệ.
    \item imageUrl: tối đa 500 ký tự.
    \item isAvailable: true/false.
\end{itemize}

\textbf{2. Read -- Xem danh sách / chi tiết sản phẩm}
\begin{itemize}
    \item Xem danh sách toàn bộ sản phẩm (phân trang).
    \item Cho phép tìm kiếm theo tên, category.
    \item Cho phép xem chi tiết theo ID.
\end{itemize}

\textbf{3. Update -- Cập nhật thông tin}
\begin{itemize}
    \item Cho phép cập nhật tất cả trừ createdBy.
    \item Validation giống Create.
    \item Không cho update nếu sản phẩm không tồn tại.
\end{itemize}

\textbf{4. Delete -- Xóa sản phẩm}
\begin{itemize}
    \item Xóa bằng ID. Nếu không tồn tại trả về "Product not found".
\end{itemize}

\subsubsection{Test Scenarios cho Product (2 điểm)}
\begin{longtable}{|p{0.05\linewidth}|p{0.35\linewidth}|p{0.15\linewidth}|p{0.4\linewidth}|}
\hline
\textbf{STT} & \textbf{Test Scenario} & \textbf{Loại} & \textbf{Kết quả mong đợi} \\
\hline
1 & Create product thành công & Happy path & Trả về ProductResponse đúng dữ liệu \\ \hline
2 & Update product thành công & Happy path & Dữ liệu được cập nhật \\ \hline
3 & Delete product thành công & Happy path & Sản phẩm bị xóa \\ \hline
4 & Read product by ID thành công & Happy path & Trả về thông tin sản phẩm \\ \hline
5 & Create thất bại (Name rỗng) & Negative & Lỗi "Tên sản phẩm không được để trống" \\ \hline
6 & Create thất bại (Price <= 0) & Negative & Lỗi "Giá phải lớn hơn 0" \\ \hline
7 & Description > 1000 ký tự & Boundary & Lỗi validation \\ \hline
8 & Update sản phẩm không tồn tại & Edge case & Lỗi "Không tìm thấy sản phẩm..." \\ \hline
9 & Delete ID không tồn tại & Edge case & Lỗi "Không tìm thấy sản phẩm..." \\ \hline
10 & Search product theo tên & Functional & Trả về danh sách phù hợp \\ \hline
\end{longtable}

\subsubsection{Thiết kế Test Cases chi tiết (5 điểm)}

% Product TC 1
\begin{longtable}{|p{0.3\linewidth}|p{0.65\linewidth}|}
\hline
\textbf{Test Case ID} & TC\_PRODUCT\_001 \\ \hline
\textbf{Test Name} & Tạo sản phẩm mới thành công \\ \hline
\textbf{Priority} & Critical \\ \hline
\textbf{Preconditions} & User đã đăng nhập, có quyền tạo sản phẩm \\ \hline
\textbf{Test Steps} & 1. Navigate to Product page \newline 2. Click "Add New Product" \newline 3. Nhập thông tin hợp lệ \newline 4. Click Save \\ \hline
\textbf{Test Data} & Name: Laptop Dell XPS, Price: 25000000, Quantity: 10, Category: Electronics, isAvailable: true \\ \hline
\textbf{Expected Result} & - Sản phẩm tạo thành công \newline - API trả về ProductResponse có ID mới \newline - Hiển thị thông báo thành công \\ \hline
\textbf{Status} & Not Run \\ \hline
\end{longtable}

% Product TC 2
\begin{longtable}{|p{0.3\linewidth}|p{0.65\linewidth}|}
\hline
\textbf{Test Case ID} & TC\_PRODUCT\_002 \\ \hline
\textbf{Test Name} & Cập nhật sản phẩm thành công \\ \hline
\textbf{Priority} & High \\ \hline
\textbf{Test Steps} & 1. Navigate to Product page \newline 2. Chọn sản phẩm ID = 5 \newline 3. Click "Edit Product" \newline 4. Nhập thông tin hợp lệ \newline 5. Click Save \\ \hline
\textbf{Expected Result} & - Cập nhật thành công \newline - API trả về ProductResponse mới nhất \\ \hline
\end{longtable}

% Product TC 3
\begin{longtable}{|p{0.3\linewidth}|p{0.65\linewidth}|}
\hline
\textbf{Test Case ID} & TC\_PRODUCT\_003 \\ \hline
\textbf{Test Name} & Xóa sản phẩm thành công \\ \hline
\textbf{Priority} & Critical \\ \hline
\textbf{Test Steps} & 1. Navigate to Product page \newline 2. Chọn sản phẩm ID = 10 \newline 3. Click "Delete Product" \newline 4. Confirm Delete \\ \hline
\textbf{Expected Result} & - Xóa thành công \newline - Sản phẩm biến mất khỏi danh sách \\ \hline
\end{longtable}

% Product TC 4
\begin{longtable}{|p{0.3\linewidth}|p{0.65\linewidth}|}
\hline
\textbf{Test Case ID} & TC\_PRODUCT\_004 \\ \hline
\textbf{Test Name} & Xem chi tiết sản phẩm theo ID \\ \hline
\textbf{Priority} & Medium \\ \hline
\textbf{Test Steps} & 1. Navigate to Product page \newline 2. Click vào sản phẩm ID = 3 \\ \hline
\textbf{Expected Result} & - API trả về ProductResponse đầy đủ các trường \\ \hline
\end{longtable}

% Product TC 5
\begin{longtable}{|p{0.3\linewidth}|p{0.65\linewidth}|}
\hline
\textbf{Test Case ID} & TC\_PRODUCT\_005 \\ \hline
\textbf{Test Name} & Tạo sản phẩm thất bại do giá không hợp lệ \\ \hline
\textbf{Priority} & High \\ \hline
\textbf{Test Steps} & 1. Navigate to Product page \newline 2. Click "Add New Product" \newline 3. Nhập Price = 0 \newline 4. Click Save \\ \hline
\textbf{Expected Result} & - Hệ thống trả lỗi validation: "Giá phải lớn hơn 0" \newline - API trả HTTP 400 \\ \hline
\end{longtable}

\section{CÂU 2: UNIT TESTING VÀ TDD}
Nhóm áp dụng quy trình TDD (Test-Driven Development) theo chu trình: \textbf{Red} (Viết test fail) $\rightarrow$ \textbf{Green} (Viết code để pass test) $\rightarrow$ \textbf{Refactor} (Tối ưu code).

\subsection{Frontend Unit Tests (TDD Approach)}

\subsubsection{Phát triển unit tests cho validation module của Login}

\textbf{Bước 1: Red (Viết Test)}
\begin{lstlisting}[language=JavaScript, caption=Login Validation Test Script (Red Phase)]
 import { validateUsername, validatePassword } from '../utils/validation';
describe('Validation Utils Tests', () => {
  // ========== a) Unit tests cho validateUsername() (2 điểm) =============
  describe('validateUsername()', () => {
    // Test username rỗng
    test('Nên trả về lỗi khi username rỗng', () => {
      expect(validateUsername('')).toBe('Username không được để trống');
      expect(validateUsername(null)).toBe('Username không được để trống');
      expect(validateUsername(undefined)).toBe('Username không được để trống');
    });
    // Test username chỉ có khoảng trắng
    test('Nên trả về lỗi khi username chỉ có khoảng trắng', () => {
      expect(validateUsername('   ')).toBe('Username không được để trống');
      expect(validateUsername('  ')).toBe('Username không được để trống');
    });
    // Test username quá ngắn
    test('Nên trả về lỗi khi username quá ngắn (< 3 ký tự)', () => {
      expect(validateUsername('a')).toBe('Username phải có ít nhất 3 ký tự');
      expect(validateUsername('ab')).toBe('Username phải có ít nhất 3 ký tự');
    });
    // Test username quá dài
    test('Nên trả về lỗi khi username quá dài (> 20 ký tự)', () => {
      const longUsername = 'a'.repeat(21);
      expect(validateUsername(longUsername)).toBe('Username không được vượt quá 20 ký tự');
      expect(validateUsername('abcdefghijklmnopqrstu')).toBe('Username không được vượt quá 20 ký tự');
    });
    // Test ký tự đặc biệt không hợp lệ
    test('Nên trả về lỗi khi username chứa ký tự đặc biệt không hợp lệ', () => {
      expect(validateUsername('user@name')).toBe('Username chỉ được chứa chữ cái, số và dấu gạch dưới');
      expect(validateUsername('user name')).toBe('Username chỉ được chứa chữ cái, số và dấu gạch dưới');
      expect(validateUsername('user-name')).toBe('Username chỉ được chứa chữ cái, số và dấu gạch dưới');
      expect(validateUsername('user#123')).toBe('Username chỉ được chứa chữ cái, số và dấu gạch dưới');
      expect(validateUsername('user.name')).toBe('Username chỉ được chứa chữ cái, số và dấu gạch dưới');
    });
    // Test username hợp lệ
    test('Nên trả về null khi username hợp lệ', () => {
      expect(validateUsername('abc')).toBeNull();
      expect(validateUsername('user123')).toBeNull();
      expect(validateUsername('test_user')).toBeNull();
      expect(validateUsername('User_123')).toBeNull();
      expect(validateUsername('abcdefghij')).toBeNull();
      expect(validateUsername('abcdefghijklmnopqrst')).toBeNull();
    });
    // Test edge cases
    test('Nên xử lý đúng các edge cases', () => {
      expect(validateUsername('___')).toBeNull(); // Chỉ gạch dưới
      expect(validateUsername('123')).toBeNull(); // Chỉ số
      expect(validateUsername('ABC')).toBeNull(); // Chữ hoa
      expect(validateUsername('abc123___')).toBeNull(); // Mix hợp lệ
    });
  });
  // ========= b) Unit tests cho validatePassword() (2 điểm) ==============
  describe('validatePassword()', () => {
    // Test password rỗng
    test('Nên trả về lỗi khi password rỗng', () => {
      expect(validatePassword('')).toBe('Password không được để trống');
      expect(validatePassword(null)).toBe('Password không được để trống');
      expect(validatePassword(undefined)).toBe('Password không được để trống');
    });
    // Test password chỉ có khoảng trắng
    test('Nên trả về lỗi khi password chỉ có khoảng trắng', () => {
      expect(validatePassword('      ')).toBe('Password không được để trống');
      expect(validatePassword('  ')).toBe('Password không được để trống');
    });
    // Test password quá ngắn
    test('Nên trả về lỗi khi password quá ngắn (< 6 ký tự)', () => {
      expect(validatePassword('12345')).toBe('Password phải có ít nhất 6 ký tự');
      expect(validatePassword('abc')).toBe('Password phải có ít nhất 6 ký tự');
      expect(validatePassword('a')).toBe('Password phải có ít nhất 6 ký tự');
      expect(validatePassword('Test1')).toBe('Password phải có ít nhất 6 ký tự');
    });
    // Test password quá dài
    test('Nên trả về lỗi khi password quá dài (> 30 ký tự)', () => {
      const longPassword = 'A1' + 'a'.repeat(30);
      expect(validatePassword(longPassword)).toBe('Password không được vượt quá 30 ký tự');
      expect(validatePassword('Test123' + 'a'.repeat(25))).toBe('Password không được vượt quá 30 ký tự');
    });
    // Test password không có chữ cái
    test('Nên trả về lỗi khi password không có chữ cái', () => {
      expect(validatePassword('123456')).toBe('Password phải chứa ít nhất một chữ cái');
      expect(validatePassword('!@#$%^')).toBe('Password phải chứa ít nhất một chữ cái');
      expect(validatePassword('12345678')).toBe('Password phải chứa ít nhất một chữ cái');
    });
    // Test password không có số
    test('Nên trả về lỗi khi password không có số', () => {
      expect(validatePassword('abcdef')).toBe('Password phải chứa ít nhất một chữ số');
      expect(validatePassword('Password')).toBe('Password phải chứa ít nhất một chữ số');
      expect(validatePassword('TestPass')).toBe('Password phải chứa ít nhất một chữ số');
    });
    // Test password hợp lệ
    test('Nên trả về null khi password hợp lệ', () => {
      expect(validatePassword('Pass123')).toBeNull();
      expect(validatePassword('Test123')).toBeNull();
      expect(validatePassword('abcdef1')).toBeNull();
      expect(validatePassword('MyP4ssw0rd')).toBeNull();
      expect(validatePassword('Test@123')).toBeNull(); // Có ký tự đặc biệt vẫn OK
      expect(validatePassword('a1b2c3')).toBeNull(); // Đúng 6 ký tự (min)
      expect(validatePassword('Test123' + 'a'.repeat(23))).toBeNull(); // 30 ký tự (max)
    });
    // Test edge cases - password với ký tự đặc biệt
    test('Nên chấp nhận password có ký tự đặc biệt (nếu có chữ và số)', () => {
      expect(validatePassword('P@ssw0rd')).toBeNull();
      expect(validatePassword('Test!123')).toBeNull();
      expect(validatePassword('My#Pass1')).toBeNull();
    });
    // Test edge cases - mix chữ hoa/thường/số
    test('Nên chấp nhận password mix chữ hoa, thường và số', () => {
      expect(validatePassword('ABCdef123')).toBeNull();
      expect(validatePassword('Test123TEST')).toBeNull();
      expect(validatePassword('aB1cD2eF3')).toBeNull();
    });
  });
  // ==================== c) Coverage tests (1 điểm) ====================
  describe('Edge Cases và Coverage', () => { 
    test('validateUsername - test với giá trị boolean/number', () => {
      // Test với các kiểu dữ liệu khác
      expect(validateUsername(123)).toBe('Username không được để trống');
      expect(validateUsername(true)).toBe('Username không được để trống');
      expect(validateUsername(false)).toBe('Username không được để trống');
    });
    test('validatePassword - test với giá trị boolean/number', () => {
      // Test với các kiểu dữ liệu khác
      expect(validatePassword(123456)).toBe('Password không được để trống');
      expect(validatePassword(true)).toBe('Password không được để trống');
      expect(validatePassword(false)).toBe('Password không được để trống');
    });
    test('validateUsername - test boundary values', () => {
      // Boundary: đúng 3 ký tự (min valid)
      expect(validateUsername('abc')).toBeNull(); 
      // Boundary: đúng 20 ký tự (max valid)
      expect(validateUsername('12345678901234567890')).toBeNull();
      // Boundary: 2 ký tự (min - 1)
      expect(validateUsername('ab')).toBe('Username phải có ít nhất 3 ký tự');     
      // Boundary: 21 ký tự (max + 1)
      expect(validateUsername('123456789012345678901')).toBe('Username không được vượt quá 20 ký tự');
    });
    test('validatePassword - test boundary values', () => {
      // Boundary: đúng 6 ký tự (min valid)
      expect(validatePassword('Pass12')).toBeNull();   
      // Boundary: đúng 30 ký tự (max valid)
      expect(validatePassword('Pass1' + 'a'.repeat(25))).toBeNull();    
      // Boundary: 5 ký tự (min - 1)
      expect(validatePassword('Pas12')).toBe('Password phải có ít nhất 6 ký tự');   
      // Boundary: 31 ký tự (max + 1)
      expect(validatePassword('Pass1' + 'a'.repeat(26))).toBe('Password không được vượt quá 30 ký tự');
    });
  });
});

\end{lstlisting}

\textbf{Bước 2: Green (Viết Code)}
\begin{lstlisting}[language=JavaScript, caption=TDD Cycle 1 - Implementation (Green)]
  if (!name || name.trim() === '') { return 'Tên sản phẩm không được để trống';  }
  // Yêu cầu: 3-100 ký tự
  if (name.length < 3 || name.length > 100) { return 'Tên sản phẩm phải từ 3-100 ký tự';  }
  return ''; };

export const validatePrice = (price) => {
  if (price === null || price === undefined || price === '') { return 'Giá không được để trống';  }
  const numPrice = Number(price);
  if (isNaN(numPrice)) {return 'Giá phải là số';  }
  if (numPrice <= 0) {return 'Giá phải lớn hơn 0'; }
  // Yêu cầu: Max 999,999,999
  if (numPrice > 999999999) { return 'Giá quá lớn (tối đa 999,999,999)'; }
  return ''; };

export const validateQuantity = (quantity) => {
  if (quantity === null || quantity === undefined || quantity === '') {return 'Số lượng không được để trống';  }
  const numQuantity = Number(quantity);
  if (isNaN(numQuantity)) { return 'Số lượng phải là số';  }
  if (!Number.isInteger(numQuantity)) { return 'Số lượng phải là số nguyên'; }
  if (numQuantity < 0) { return 'Số lượng không được âm'; }
  if (numQuantity > 99999) {return 'Số lượng không được quá 99,999'; }
  return ''; };

export const validateDescription = (description) => {
  if (description && description.length > 500) {return 'Mô tả không được quá 500 ký tự';  }
  return '';};
export const validateCategory = (category) => {
  const VALID_CATEGORIES = ['Electronics', 'Fashion', 'Food', 'Books'];
    if (!category || category.trim() === '') {    return 'Danh mục không được để trống';  }
    if (!VALID_CATEGORIES.includes(category)) { return 'Danh mục không hợp lệ';  }
  return ''; };
/**
 * Format currency (VND)
 */
export const formatCurrency = (amount) => {
  return new Intl.NumberFormat('vi-VN', {
    style: 'currency',
    currency: 'VND'
  }).format(amount);
};
/**
 * Format date
 */
export const formatDate = (dateString) => {
  const date = new Date(dateString);
  return date.toLocaleDateString('vi-VN');
};

\end{lstlisting}
\begin{figure}[htbp]
  \centering
  \includegraphics[width=0.8\textwidth]{1.jpg}
  \caption{Mô tả ảnh}
  \label{fig:tdd-refactor-1}
\end{figure}
\textbf{Bước 3:Refactor (Tối ưu)}
Do logic validation product khá nhiều điều kiện if-else, trong bước Refactor, nhóm tiến hành kiểm tra lại coverage để đảm bảo các nhánh (branch) điều kiện như giá min/max, số lượng min/max đều được bao phủ.

\subsubsection{Phát triển unit tests cho validation module của Product}
Tương tự Login, module Product cũng phát triển theo TDD.

\textbf{Bước 1: Red (Viết Test Case)}
\begin{lstlisting}[language=JavaScript, caption=Product Validation Test Script]
import { validateProductName, validatePrice, validateQuantity } from '../utils/validation.js';

describe('Product Validation Utilities Tests', () => {
    describe('validateProductName', () => {
        test('returns error for empty product name', () => {
            expect(validateProductName('')).toBe('Tên sản phẩm không được để trống');
        });
        test('returns error for short product name (< 3 chars)', () => {
            expect(validateProductName('AB')).toBe('Tên sản phẩm phải từ 3-100 ký tự');
        });
    });

    describe('validatePrice', () => {
        test('returns error for empty price', () => {
            expect(validatePrice('')).toBe('Giá không được để trống');
        });
        test('returns error for price <= 0', () => {
            expect(validatePrice(0)).toBe('Giá phải lớn hơn 0');
        });
    });
    
    // ... validateQuantity tests ...
});
\end{lstlisting}

\textbf{Bước 2: Green (Viết Code Logic)}
\begin{lstlisting}[language=JavaScript, caption=Product Implementation]
export const validateProductName = (name) => {
    if (!name || name.trim() === '') {
        return 'Tên sản phẩm không được để trống';
    }
    if (name.length < 3 || name.length > 100) {
        return 'Tên sản phẩm phải từ 3-100 ký tự';
    }
    return '';
};

export const validatePrice = (price) => {
    if (price === null || price === undefined || price === '') {
        return 'Giá không được để trống';
    }
    const numPrice = Number(price);
    if (isNaN(numPrice)) { return 'Giá phải là số'; }
    if (numPrice <= 0) { return 'Giá phải lớn hơn 0'; }
    return '';
};
\end{lstlisting}

\subsection{Backend Unit Tests (JUnit 5)}

\subsubsection{Login Service Tests}
\begin{lstlisting}[language=Java, caption=AuthService Test]
@ExtendWith(MockitoExtension.class)
@DisplayName("TC1: Login thanh cong voi credentials hop le")
public class AuthServiceTest {
    @Mock private UserRepository userRepository;
    @Mock private PasswordEncoder passwordEncoder;
    @Mock private AuthenticationManager authenticationManager;
    @Mock private JwtUtils jwtUtils;
    @Mock private Authentication authentication;
    @InjectMocks private AuthService authService;

    // ... setup vars ...

    @Test
    void testLogin_Success() {
        when(authenticationManager.authenticate(any(UsernamePasswordAuthenticationToken.class)))
            .thenReturn(authentication);
        when(userRepository.findByUsername(anyString())).thenReturn(Optional.of(user));
        
        AuthResponse response = authService.login(loginRequest);
        assertNotNull(response);
        assertEquals("fake-jwt-token", response.getToken());
        verify(jwtUtils).generateToken("testuser");
    }

    @Test
    @DisplayName("TC2: Login that bai voi username khong ton tai")
    void testLogin_UserNotFound() {
        when(userRepository.findByUsername(anyString())).thenReturn(Optional.empty());
        RuntimeException exception = assertThrows(RuntimeException.class, () -> { 
            authService.login(loginRequest);
        });
        assertEquals("User không tồn tại!", exception.getMessage());
    }
}
\end{lstlisting}

\subsubsection{Product Service Tests}
\begin{lstlisting}[language=Java, caption=ProductService Test with Fake Repository]
@DisplayName("Product Service Unit tests")
public class ProductServiceTest {
    private ProductService productService;
    private ProductRepository productRepository; // Mocked manually in source
    // ... setup fake lists ...
    
    @Test
    @DisplayName("TC1: Tao san pham moi thanh cong")
    void testCreateProduct() {
        ProductRequestDTO request = new ProductRequestDTO("Asus", BigDecimal.valueOf(15000000), 10, "Desc", 1L);
        ProductResponseDTO result = productService.createProduct(request);
        assertNotNull(result);
        assertEquals("Asus", result.getProductName());
    }
    
    // ... other tests (update, delete, get) ...
}
\end{lstlisting}

\section{CÂU 3: INTEGRATION TESTING}

\subsection{Login - Integration Testing}
\subsubsection{Test tích hợp Login component với API service}
\paragraph{a) Test rendering \& user interactions} \mbox{}\\
\begin{lstlisting}[language=JavaScript, caption=Rendering test]
test('Rendering component và user interactions cơ bản', () => {
  render(
    <BrowserRouter future={{ v7_startTransition: true, v7_relativeSplatPath: true }}>
      <Login />
    </BrowserRouter>
  );
  // Kiểm tra rendering đầy đủ - dùng getByRole để tránh duplicate
  expect(screen.getByTestId('username-input')).toBeInTheDocument();
  expect(screen.getByTestId('password-input')).toBeInTheDocument();
  expect(screen.getByTestId('login-button')).toBeInTheDocument();
  expect(screen.getByRole('heading', { name: /Đăng Nhập/i })).toBeInTheDocument(); // Chỉ lấy h2
  expect(screen.getByText(/Chào mừng bạn trở lại/i)).toBeInTheDocument();
  // Test user interactions - username
  const usernameInput = screen.getByTestId('username-input');
  fireEvent.change(usernameInput, { target: { value: 'testuser' } });
  expect(usernameInput.value).toBe('testuser');
  // Test user interactions - password
  const passwordInput = screen.getByTestId('password-input');
  fireEvent.change(passwordInput, { target: { value: 'Test123' } });
  expect(passwordInput.value).toBe('Test123');  });


\end{lstlisting}
\paragraph{b) Test form submission & API calls} \mbox{}\\
\begin{lstlisting}[language=JavaScript, caption=Gọi API khi submit form hợp lệ]
test('Gọi API khi submit form hợp lệ và handling success', async () => {
    const mockNavigate = jest.fn();
    const mockOnLogin = jest.fn();
    useNavigate.mockReturnValue(mockNavigate);
    authService.login.mockResolvedValue({
      success: true,
      message: 'Đăng nhập thành công!',
      data: { token: 'fake-token' }
    });
    render(
      <BrowserRouter future={{ v7_startTransition: true, v7_relativeSplatPath: true }}>
        <Login onLogin={mockOnLogin} />
      </BrowserRouter>    );
    fireEvent.change(screen.getByTestId('username-input'), {
      target: { value: 'testuser' }
    });
    fireEvent.change(screen.getByTestId('password-input'), {
      target: { value: 'Test123' }
    });
    fireEvent.click(screen.getByTestId('login-button'));
    await waitFor(() => { expect(authService.login).toHaveBeenCalledWith('testuser', 'Test123');  });
    await waitFor(() => { expect(mockOnLogin).toHaveBeenCalled();  });
    await waitFor(() => { expect(mockNavigate).toHaveBeenCalledWith('/products'); });
  });

\end{lstlisting}
\paragraph{c) Test Error Handling & Success Messages} \mbox{}\\
\begin{lstlisting}[language=JavaScript, caption= Hiển thị thông báo thành công nếu API trả về thành công]
test('Gọi API khi submit form hợp lệ và handling success', async () => {
    const mockNavigate = jest.fn();
    const mockOnLogin = jest.fn();
    useNavigate.mockReturnValue(mockNavigate);
    authService.login.mockResolvedValue({
      success: true,
      message: 'Đăng nhập thành công!',
      data: { token: 'fake-token' }  });
    render(
      <BrowserRouter future={{ v7_startTransition: true, v7_relativeSplatPath: true }}>
        <Login onLogin={mockOnLogin} />
      </BrowserRouter>  );
    fireEvent.change(screen.getByTestId('username-input'), {
      target: { value: 'testuser' } });
    fireEvent.change(screen.getByTestId('password-input'), {
      target: { value: 'Test123' }  });
    fireEvent.click(screen.getByTestId('login-button'));
    await waitFor(() => { expect(authService.login).toHaveBeenCalledWith('testuser', 'Test123'); });
    await waitFor(() => { expect(mockOnLogin).toHaveBeenCalled(); });
    await waitFor(() => { expect(mockNavigate).toHaveBeenCalledWith('/products'); });
  });


\end{lstlisting}
\begin{lstlisting}[language=JavaScript, caption= Hiển thị thông báo thành công nếu API trả về Lỗi]
    test('Handling error khi API fail', async () => {
    authService.login.mockRejectedValue({ response: { data: { message: 'Lỗi server' } }  });
    render(
      <BrowserRouter future={{ v7_startTransition: true, v7_relativeSplatPath: true }}>
        <Login />
      </BrowserRouter> );
    fireEvent.change(screen.getByTestId('username-input'), { target: { value: 'testuser' } });
    fireEvent.change(screen.getByTestId('password-input'), { target: { value: 'Test123' } });
    fireEvent.click(screen.getByTestId('login-button'));
    expect(await screen.findByText('Lỗi server')).toBeInTheDocument(); });

test('Hiển thị default error message khi API fail không có message', async () => {
    authService.login.mockRejectedValue(new Error('Network error'));
    render(
      <BrowserRouter future={{ v7_startTransition: true, v7_relativeSplatPath: true }}>
        <Login />
      </BrowserRouter> );
    fireEvent.change(screen.getByTestId('username-input'), { target: { value: 'testuser' }});
    fireEvent.change(screen.getByTestId('password-input'), { target: { value: 'Test123' }});
    fireEvent.click(screen.getByTestId('login-button'));
    expect(await screen.findByText('Tên đăng nhập hoặc mật khẩu không đúng!'))
.toBeInTheDocument();
  });

\end{lstlisting}

\subsubsection{Test API endpoints của Login với MockMvc}
\paragraph{a) Test POST /api/auth/login endpoint & Test response structure và status codes}\mbox{}
\begin{lstlisting}[language=Java, caption=POST /api/auth/login endpoint - SUCCESS]
 @Test
 @DisplayName("POST /api/auth/login endpoint - SUCCESS")
 void testLoginSuccess() throws Exception {
	LoginRequest request = new LoginRequest(
	"testuser",
	"Test123"  );
	AuthResponse mockResponse = new AuthResponse(
	"token123", 1L,  "testuser", "test@gmail.com", "USER"  );
	when(authService.login(any(LoginRequest.class)))
	.thenReturn(mockResponse);
	mockMvc.perform(post("/api/auth/login")
	.contentType(MediaType.APPLICATION_JSON)
	.content(objectMapper.writeValueAsString(request )))
	.andExpect(status().isOk())
	.andExpect(jsonPath("$.success").value(true))
	.andExpect(jsonPath("$.message").value("Đăng nhập thành công!"))
	.andExpect(jsonPath("$.data.token").value("token123"))
	.andExpect(jsonPath("$.data.id").value(1L))
	.andExpect(jsonPath("$.data.username").value("test@gmail.com"))
	.andExpect(jsonPath("$.data.email").value(1L))
	.andExpect(jsonPath("$.data.role").value("USER"))
	;
 }

\end{lstlisting}

\begin{lstlisting}[language=javaScript, caption=POST /api/auth/login endpoint - FAILED username không tồn tại]
    Copy code dáng vào đây
\end{lstlisting}

\begin{lstlisting}[language=javaScript, caption=POST /api/auth/login endpoint - FAILED Password sai]
    Copy code dáng vào đây
\end{lstlisting}
\paragraph{c) Test CORS và headers}\mbox{}\\
\begin{lstlisting}[language=javaScript, caption= Kiểm tra CORS headers cho POST request
]
    Copy code dáng vào đây
\end{lstlisting}
\subsection{Product - Integration Testing}
\subsubsection{Test Product component với API service}
\paragraph { a) Productlist compoment \& API} \mbox{}\\
\begin{lstlisting}[language=JavaScript, caption=ProductList Integration]

  test('ProductList tải và hiển thị danh sách sản phẩm', async () => {
    productService.getAllProducts.mockResolvedValue({
      success: true,
      data: mockProducts
    });

    render(
      <MemoryRouter>
        <ProductList onLogout={jest.fn()} />
      </MemoryRouter>
    );

    expect(await screen.findByText(/Product Management/i)).toBeInTheDocument();

    await waitFor(
      () => {
        expect(
          screen.queryByText(/Đang tải sản phẩm/i)
        ).not.toBeInTheDocument();
      },
      { timeout: 5000 }
    );

    // Kiểm tra sản phẩm được hiển thị
    expect(await screen.findByText(/iPhone 15 Pro/i)).toBeInTheDocument();
    expect(screen.getByText(/Galaxy S24 Ultra/i)).toBeInTheDocument();
  });
\end{lstlisting}
\begin{lstlisting} [languege=javaScript,caption=Productlist empty]
    test('ProductList hiển thị empty state khi không có sản phẩm', async () => {
   productService.getAllProducts.mockResolvedValue({
      success: true,
      data: []
    });
    render(
      <MemoryRouter>
        <ProductList onLogout={jest.fn()} />
      </MemoryRouter>
    );
    await waitFor(() => {
      expect(screen.getByText(/Chưa có sản phẩm nào/i)).toBeInTheDocument();
    });
  });
\end{lstlisting}
\paragraph{b) ProductForm \& API }\mbox{}\\
\begin{lstlisting}[language=javaScript, caption= Form thêm sản phẩm
]
    Copy code dáng vào đây
\end{lstlisting}
\begin{lstlisting}[language=javaScript, caption= Form sửa sản phẩm
]
    Copy code dáng vào đây
\end{lstlisting}
\paragraph{b) ProductDetail \& API }\mbox{}\\
\begin{lstlisting}[language=javaScript, caption= Form Product detail
]
    Copy code dáng vào đây
\end{lstlisting}
\begin{lstlisting}[language=javaScript, caption= Form sửa sản phẩm
]
    Copy code dáng vào đây
\end{lstlisting}

\subsubsection{Test các API endpoints của Product}
\paragraph{a) Create product - POST /api/products \& API }\mbox{}\\
\begin{lstlisting}[language=javaScript, caption= Viết mô tả code
]
    Copy code dáng vào đây
\end{lstlisting}
\paragraph{b) Get all product - GET /api/products\& API }\mbox{}\\
\begin{lstlisting}[language=javaScript, caption= Viết mô tả code
]
    Copy code dáng vào đây
\end{lstlisting}

\paragraph{c) Test GET PRODUCT - GET /api/products/{id}\& API }\mbox{}\\
\begin{lstlisting}[language=javaScript, caption= Viết mô tả code
]
    Copy code dáng vào đây
\end{lstlisting}

\paragraph{d) 
Test UPDATE PRODUCT - PUT /api/products/{id} }\mbox{}\\
\begin{lstlisting}[language=javaScript, caption= Viết mô tả code
]
    Copy code dáng vào đây
\end{lstlisting}

Test UPDATE PRODUCT - DELETE /api/products/{id} }\mbox{}\\
\begin{lstlisting}[language=javaScript, caption= Viết mô tả code
]
    Copy code dáng vào đây
\end{lstlisting}

\section{CÂU 4: MOCK TESTING}

\subsection{Login - Mock Testing}
\subsubsection{Frontend Mocking}
\paragraph{Mục tiêu test:}Mock external dependencies cho Login component nhằm kiểm tra logic và hành vi của component Login trong các tình huống khác nhau mà không phụ thuộc vào backend thực tế.
\paragraph{Các bước thực hiện}\mbox{}\\
\begin{itemize}
  \item Mock \texttt{authService.loginUser()}
  \item Test với mocked successful/failed responses
  \item Verify mock calls
\end{itemize}

\paragraph{Các test case đã thực hiện}\mbox{}\\

\begin{enumerate}
  \item Kiểm tra rendering và tương tác người dùng (2 điểm): Hiển thị form và cho phép nhập dữ liệu
  \item Kiểm tra form submission với mocked successful response
  \item Kiểm tra xử lý lỗi với mocked failed response
  \item Verify mock calls – kiểm tra mock được reset
\end{enumerate}
\begin{lstlisting}[language=JavaScript, caption=Mock Login Service]
    copy code vào đây
\end{lstlisting}

\begin{lstlisting}[language=JavaScript, caption=Thêm mô tả]
    copy code vào đây
\end{lstlisting}
\begin{lstlisting}[language=JavaScript, caption=Viết mô tả và đây]
    copy code vào đây
    muốn thêm code block thì tự copy pass
\end{lstlisting}
\subsubsection{Backend Mocking}
\paragraph{Mục tiêu test:}
Mock dependencies trong Backend tests để kiểm tra logic của controller mà không phụ thuộc vào service thực tế, đảm bảo controller xử lý đúng các request và response.

\paragraph{Các bước thực hiện:}
\begin{itemize}
  \item Mock \texttt{AuthService} với \texttt{@MockBean}
  \item Test controller với mocked service
  \item Verify mock interactions
\end{itemize}

\paragraph{Các test case đã thực hiện:}
\begin{enumerate}[label=TC\arabic*:]
  \item Login thành công với mocked service
  \item Login thất bại -- username rỗng
  \item Login thất bại -- password rỗng
  \item Login thất bại -- password không có số
  \item Login thất bại -- username quá dài
\end{enumerate}

\paragraph{Kết quả test:}\mbox{}\\
\begin{lstlisting}[language=Java, caption=Mô tả test]
    Copy code vào đây muốn thêm thì tự copy pass
}
\end{lstlisting}

\subsection{Product - Mock Testing}
\subsubsection{Frontend Mocking}
\paragraph{Mục tiêu test:}
Mock \texttt{ProductService} trong component tests để kiểm tra các operations CRUD mà không phụ thuộc vào service thực tế.

\paragraph{Các bước thực hiện:}
\begin{itemize}
  \item Mock các thao tác CRUD của \texttt{ProductService}
  \item Test các kịch bản thành công và thất bại
  \item Verify tất cả các mock calls
\end{itemize}

\paragraph{Các test case đã thực hiện:}
\begin{enumerate}
  \item Mock: Get products success – Kiểm tra việc lấy danh sách sản phẩm thành công
  \item Mock: Get products failure – Kiểm tra việc lấy danh sách sản phẩm thất bại
  \item Mock: Create product success – Kiểm tra việc tạo sản phẩm mới thành công
  \item Mock: Create product failure – Kiểm tra việc tạo sản phẩm mới thất bại
  \item Mock: Update product success – Kiểm tra việc cập nhật sản phẩm thành công
  \item Mock: Update product failure – Kiểm tra việc cập nhật sản phẩm thất bại
  \item Mock: Delete product success – Kiểm tra việc xóa sản phẩm thành công
\end{enumerate}

\paragraph{Kết quả test:}\mbox{}\\
\begin{figure}[htbp]
  \centering
  \includegraphics[width=0.8\textwidth]{521.png}
  \caption{Mô tả ảnh}
  \label{fig:test-result-521}
\end{figure}
\begin{lstlisting}[language=JavaScript, caption=Chi tiết test case]
import React from 'react';
import { render, screen, waitFor } from '@testing-library/react';
import ProductList from '../components/ProductList.js';
import { MemoryRouter } from 'react-router-dom';

// --- SỬA QUAN TRỌNG: Bỏ "* as" đi ---
import productService from '../services/productService.js'; 

// Mock module: Định nghĩa rõ cấu trúc để Jest hiểu "default export" có những hàm nào
jest.mock('../services/productService.js', () => ({
  __esModule: true, // Báo cho Jest biết đây là ES Module
  default: {
    getAllProducts: jest.fn(),
    createProduct: jest.fn(),
    updateProduct: jest.fn(),
    deleteProduct: jest.fn(),
    getProductById: jest.fn(),  }, }));
const mockProducts = [
  { id: 1, name: 'Iphone', price: 10000000 },
  { id: 2, name: 'Samsung', price: 8000000 } ];
describe('Product Service Mock CRUD (No UI Interaction)', () => {  
  beforeEach(() => {
    jest.clearAllMocks();  });
  // --- 1. GET ALL (Integration nhẹ: component gọi service khi render) ---
  test('getAllProducts - success: Render component calls API', async () => {
    // Setup Mock
    productService.getAllProducts.mockResolvedValue({ success: true, data: mockProducts });

    render(
      <MemoryRouter>
        <ProductList />
      </MemoryRouter>    );

    // Verify Call
    await waitFor(() => {
      expect(productService.getAllProducts).toHaveBeenCalledTimes(1);    });
        // (Optional) Check UI để chắc chắn component nhận data
    // expect(screen.getByText(/Iphone/i)).toBeInTheDocument();  });
  test('getAllProducts - failure: Show error message', async () => {
    productService.getAllProducts.mockRejectedValue(new Error('Network error'));
    render(
      <MemoryRouter>
        <ProductList />
      </MemoryRouter>    );
    // Verify Call
    await waitFor(() => {
      expect(productService.getAllProducts).toHaveBeenCalledTimes(1);    });
  });

  //'2. CREATE / UPDATE / DELETE (Manual Trigger)

  // --- 1. CREATE PRODUCT ---

  test('Create product success - Kiểm tra việc tạo sản phẩm mới thành công', async () => {
    // Arrange: Chuẩn bị dữ liệu và mock trả về thành công
    const newProduct = { name: 'Samsung Galaxy S24', price: 20000000 };
    const mockResponse = { success: true, data: { id: 1, ...newProduct } };    
    productService.createProduct.mockResolvedValue(mockResponse);
    const result = await productService.createProduct(newProduct);
    expect(result).toEqual(mockResponse); // Kiểm tra dữ liệu trả về đúng
    expect(productService.createProduct).toHaveBeenCalledTimes(1);
    expect(productService.createProduct).toHaveBeenCalledWith(newProduct);  });

  test('Create product failure - Kiểm tra việc tạo sản phẩm mới thất bại', async () => {
    const newProduct = { name: 'Invalid Product' };
    const errorMessage = 'Network Error';    
    productService.createProduct.mockRejectedValue(new Error(errorMessage));
    await expect(productService.createProduct(newProduct)).rejects.toThrow(errorMessage);    
    expect(productService.createProduct).toHaveBeenCalledTimes(1);
    expect(productService.createProduct).toHaveBeenCalledWith(newProduct);  });

  // --- 2. UPDATE PRODUCT ---

  test('Update product success - Kiểm tra việc cập nhật sản phẩm thành công', async () => {
    const productId = 1;
    const updateData = { price: 18000000 };
    const mockResponse = { success: true, message: 'Updated successfully' };
    productService.updateProduct.mockResolvedValue(mockResponse);
    const result = await productService.updateProduct(productId, updateData);
    expect(result).toEqual(mockResponse);
    expect(productService.updateProduct).toHaveBeenCalledTimes(1);
    expect(productService.updateProduct).toHaveBeenCalledWith(productId, updateData);
  });

  test('Update product failure - Kiểm tra việc cập nhật sản phẩm thất bại', async () => {
    const productId = 99; // ID không tồn tại
    const updateData = { name: 'Ghost Product' };
    const errorObj = { message: 'Product not found' };
    productService.updateProduct.mockRejectedValue(errorObj);

    // Act & Assert: Sử dụng rejects.toEqual nếu lỗi trả về là object, toThrow nếu là Error instance
    await expect(productService.updateProduct(productId, updateData)).rejects.toEqual(errorObj);
    expect(productService.updateProduct).toHaveBeenCalledTimes(1);
    expect(productService.updateProduct).toHaveBeenCalledWith(productId, updateData);
  });

  // --- 3. DELETE PRODUCT ---

  test('Delete product success - Kiểm tra việc xóa sản phẩm thành công', async () => {
    const productId = 5;
    const mockResponse = { success: true };
    productService.deleteProduct.mockResolvedValue(mockResponse);
    const result = await productService.deleteProduct(productId);
    expect(result).toEqual(mockResponse);
    expect(productService.deleteProduct).toHaveBeenCalledTimes(1);
    expect(productService.deleteProduct).toHaveBeenCalledWith(productId);
  });

  test('Delete product failure - Kiểm tra việc xóa sản phẩm thất bại', async () => {
    // Arrange: Giả lập lỗi (ví dụ: lỗi mạng hoặc ID không tồn tại)
    const productId = 9999;
    const errorMessage = 'Delete failed: Product not found';
    productService.deleteProduct.mockRejectedValue(new Error(errorMessage));
    await expect(productService.deleteProduct(productId)).rejects.toThrow(errorMessage);
    expect(productService.deleteProduct).toHaveBeenCalledTimes(1);
    expect(productService.deleteProduct).toHaveBeenCalledWith(productId);
  });
});

\end{lstlisting}

\subsubsection{Backend ProductRepository Mocking}
\begin{lstlisting}[language=Java, caption=ProductService Mock Repository]
@Test
void testCreateProductSuccess() {
    ProductRequestDTO request = new ProductRequestDTO("Laptop DELL", new BigDecimal("2400"), 10, "Desc", 1L);
    // ... mock repository findById and save ...
    ProductResponseDTO result = productService.createProduct(request);
    assertNotNull(result);
    verify(productRepository, times(1)).save(any(Product.class));
}
\end{lstlisting}

\section{CÂU 5: AUTOMATION TESTING VÀ CI/CD}

\subsection{E2E Testing với Cypress}
\begin{lstlisting}[language=JavaScript, caption=Cypress E2E Login]
describe('Login E2E Tests', () => {
    beforeEach(() => { cy.visit('http://localhost:3000/login'); });

    it('Nên login thành công với credentials hợp lệ', () => {
        cy.get('[data-testid="username-input"]').type('testuser');
        cy.get('[data-testid="password-input"]').type('Test123');
        cy.get('[data-testid="login-button"]').click();
        cy.url().should('include', '/products');
    });

    it('Nên hiển thị lỗi khi submit form rỗng', () => {
        cy.get('[data-testid="login-button"]').click();
        cy.contains('Username không được để trống').should('be.visible');
    });
});
\end{lstlisting}

\subsection{CI/CD với GitHub Actions}
\begin{lstlisting}[language=yaml, caption=GitHub Actions Workflow]
name: Complete CI/CD Pipeline
on:
  push:
    branches: [ main ]
jobs:
  build:
    runs-on: ubuntu-latest
    steps:
      - uses: actions/checkout@v2
      - name: Setup Java
        uses: actions/setup-java@v2
        with: { java-version: '21' }
      - name: Backend Tests
        run: cd backend && ./mvnw clean test
      - name: Frontend Tests
        run: cd frontend && npm install && npm test
\end{lstlisting}

\section{PHẦN MỞ RỘNG (BONUS)}

\subsection{Performance Testing}
\subsubsection{k6 Script cho Login API}
\begin{lstlisting}[language=JavaScript, caption=k6 Login Script]
import http from 'k6/http';
import { check, sleep } from 'k6';

export const options = {
    stages: [
        { duration: '30s', target: 100 },
        { duration: '1m', target: 500 },
        { duration: '30s', target: 0 },
    ],
    thresholds: { http_req_duration: ['p(95)<2000'] }
};

export default function () {
    const url = 'http://localhost:8080/auth/login';
    const payload = JSON.stringify({ username: 'testuser', password: '123' });
    const params = { headers: { 'Content-Type': 'application/json' } };
    const res = http.post(url, payload, params);
    check(res, { 'status is 200': (r) => r.status === 200 });
    sleep(1);
}
\end{lstlisting}

\subsection{Security Testing}
\textbf{SQL Injection:} Sử dụng Parameterized Query (Spring Data JPA).
\begin{lstlisting}[language=Java]
@Query("SELECT u FROM User u WHERE u.username = :username")
Optional<User> findByUsername(@Param("username") String username);
\end{lstlisting}

\textbf{XSS:} React tự động escape giá trị. Backend validate input.

\section{KẾT LUẬN}
Qua bài tập lớn này, nhóm đã hoàn thành việc xây dựng và kiểm thử ứng dụng quản lý sản phẩm theo đúng quy trình:
\begin{enumerate}
    \item Phân tích yêu cầu kỹ lưỡng trước khi code.
    \item Áp dụng TDD (Test Driven Development).
    \item Thực hiện Integration Test để đảm bảo các module hoạt động tốt với nhau.
    \item Thiết lập Mock Testing để cô lập lỗi.
    \item Xây dựng pipeline CI/CD để tự động hóa quy trình kiểm thử.
\end{enumerate}
Kết quả đạt được là một hệ thống ổn định, độ bao phủ code (Code Coverage) đạt trên 80\%, đáp ứng tốt các yêu cầu đề ra của môn học.

\subsection{Tài liệu tham khảo}
\begin{enumerate}
    \item React Documentation
    \item Spring Boot Documentation
    \item Jest Documentation
    \item JUnit 5 User Guide
    \item Mockito Framework
    \item Cypress Documentation
\end{enumerate}

\end{document}